\documentclass[11pt, a4paper]{report}

\usepackage[utf8]{inputenc}
\usepackage[T1]{fontenc}
\usepackage[french]{babel}

% Appel des packages dans le dossier preambule
\def\packagepath{../../../../../preambule}   % path du package princial

\usepackage{\packagepath/page_layout} % <--- Nouveau
\usepackage{\packagepath/config_info}
\usepackage{\packagepath/cadres_cours}
\usepackage{\packagepath/zones_reponse}
\usepackage{\packagepath/config_ds}
\usepackage{\packagepath/config_maths}

% --- PILOTAGE ---
%\def\visibleornot{visible}    % visible or invisible
\ifdefined\visibleornot\else\def\visibleornot{visible}\fi


\def\documentpath{./Sources_Latex}
\graphicspath{{./Images/}}


\def\ptsexoA{0}



\begin{document}

\def\level{Premiere}  % Classe à droite
\def\course{NSI}
\def\chaptername{La boucle FOR} % Titre au centre
\def\docname{LP04~--~TD 1}        % Identifiant à gauche

\pagestyle{DS_LP}

\NomPrenom{}

\begin{exercice_cours}[pour chaque boucle, indiquer combien de fois elle s'exécute et les valeurs prises par \texttt{i}:]
    
\begin{tabular}{p{7cm}cc}
    \hline
    Instruction & Nbre de boucles & Valeurs de \texttt{i} \\
    \hline
    \verb|for i in range(5):| & \reponseLigne[3.5cm]{\texttt{5}} & \reponseLigne[4.5cm]{\texttt{0, 1, 2, 3, 4}} \\
    \verb|for i in range(2, 6):| & \reponseLigne[3.5cm]{\texttt{4}} & \reponseLigne[4.5cm]{\texttt{2, 3, 4, 5}} \\
    \verb|for i in range(0, 7):| & \reponseLigne[3.5cm]{\texttt{7}} & \reponseLigne[4.5cm]{\texttt{0, 1, 2, 3, 4, 5, 6}} \\
    \verb|for i in range(1, 10, 2):| & \reponseLigne[3.5cm]{\texttt{5}} & \reponseLigne[4.5cm]{\texttt{1, 3, 5, 7, 9}} \\
    \hline
        \end{tabular}
    \end{exercice_cours}


    
    \begin{minipage}{0.48\linewidth}
        \begin{exercice_cours}[Ecrire un code]
        Ecrire une boucle \texttt{for} qui affiche les nombres pairs de 0 à 20 (inclus).    
    
    \end{exercice_cours}

    \end{minipage}\hfill
    \begin{minipage}{0.48\linewidth}
\begin{blocvisible}[height=3.5]
    \begin{CodePythontexAlt}[TaillePolice=\large, Largeur=1\linewidth,Centre,Lignes=true,PremLigne=1]{}
def afficher_pairs():
    for i in range(0, 21, 2):
        print(i)

    \end{CodePythontexAlt}
\end{blocvisible}
    \end{minipage}

    
    \begin{minipage}{0.48\linewidth}
        \begin{exercice_cours}[Test d'appartenance]

Ecrire la fonction  qui prend en paramètre un mot et une lettre, et qui retourne: 
\begin{itemize}
        \item \texttt{True} si la lettre est présente dans le mot
        \item \texttt{False} sinon 
    \end{itemize}

    Fonction: \verb|est_present(mot, lettre_c)|
    \end{exercice_cours}

    \end{minipage}\hfill
    \begin{minipage}{0.48\linewidth}
\begin{blocvisible}[height=4.5]
    \begin{CodePythontexAlt}[TaillePolice=\large, Largeur=1\linewidth,Centre,Lignes=true,PremLigne=1]{}
def est_present(mot, lettre_c):
    for lettre in mot:
        if lettre == lettre_c:
            return True
    return False
    \end{CodePythontexAlt}
\end{blocvisible}


    \end{minipage}

        
    \begin{minipage}{0.48\linewidth}
           \begin{exercice_cours}[Nombre d'occurences]

            Ecrire la fonction \verb|compter(mot, lettre_c)| qui prend en paramètre un mot et une lettre, et qui retourne le nombre d'occurences de la lettre dans le mot.
    \end{exercice_cours}

    \end{minipage}\hfill
    \begin{minipage}{0.48\linewidth}
\begin{blocvisible}[height=4.5]
    \begin{CodePythontexAlt}[TaillePolice=\large, Largeur=1\linewidth,Centre,Lignes=true,PremLigne=1]{}
def compter(mot, lettre_c):
    count = 0
    for lettre in mot:
        if lettre == lettre_c:
            count += 1
    return count
    \end{CodePythontexAlt}
\end{blocvisible}

    \end{minipage}

\begin{minipage}{0.45\linewidth}
           \begin{exercice_cours}[Modification de str]

Ecrire la fonction \verb|remplacer(mot, L1, L2)| qui prend en paramètre un mot et deux lettres, et qui retourne le mot avec toutes les occurrences de la première lettre remplacées par la deuxième lettre.
    \end{exercice_cours}

    \end{minipage}\hfill
    \begin{minipage}{0.53\linewidth}
\begin{blocvisible}[height=5]
    \begin{CodePythontexAlt}[TaillePolice=\large, Largeur=1\linewidth,Centre,Lignes=true,PremLigne=1]{}
def remplacer(mot, L1, L2):
    nouveau_mot = ""
    for lettre in mot:
        if lettre == L1:
            nouveau_mot += L2
        else:
            nouveau_mot += lettre
    return nouveau_mot
    \end{CodePythontexAlt}
\end{blocvisible}

    \end{minipage}

\end{document}